\begin{textsample}{POS Dim 5 – global\_north\_2024\_01 – Score 19.00 – global\_north\_2024\_01/105010038.txt}  \label{ex:f5_pos_119}
Zahid Hussain ISLAMABAD -- More than two million besieged people are facing the catastrophe of famine with barely any food trickling in.
But the \textbf{international} community has failed to stop the worst \textbf{genocide} in recent history.
Israel unabatedly continues with ethnic cleansing in the enclave.
While the United States is fully backing the extermination of the occupied population the silence of the Arab countries over Israel ' s war \textbf{crimes} is deafening.
Many of them are not even willing to endorse a \textbf{case} filed by South Africa against Israel at the \textbf{International} \textbf{Court} of Justice ( ICJ ), accusing the Jewish state of \textbf{crimes} of \textbf{genocide} against Palestinians.
Their shameful capitulation seems to have encouraged Israel not only to continue its military operation on Gaza but it has also taken the war to Lebanon and other surrounding countries.
Israel ' s prime minister has made his plan clear to push out the entire population from their homes and to resettle Palestinians outside Gaza.
The Israeli government has ignored the \textbf{resolution} passed by more than 150 nations at the UN General Assembly calling for a ceasefire.
The South @ @ @ @ @ @ @ @ @ @ ' s \textbf{genocidal} actions.
It ' s apparent that American patronage and the inaction of the Arab countries has given the Jewish state complete impunity.
The latest Israeli military action inside Lebanon has already widened the war beyond Gaza and the West Bank.
The threat of regional conflagration is looming large with the increasing American military presence in the Middle East in aid of Israel.
Ironically, it has not been any Arab or Muslim country that has taken up the \textbf{genocide} \textbf{case} to the ICJ but it is South Africa that has challenged the Israeli atrocities.
The top UN \textbf{court} will hear this week an application from South Africa alleging that Israel is breaching the 1948 \textbf{Genocide} \textbf{Convention} and seeking \textbf{measures} including the immediate suspension of its military operations in Gaza.
Israel has enjoyed impunity for its \textbf{crimes} against Palestinians for almost four months.
But this situation seems to have changed after South Africa on Dec 29, 2023, \textbf{submitted} an 84-page charge sheet against the Jewish state at the ICJ.
A well-documented petition filed by the country ' s @ @ @ @ @ @ @ @ @ @ continuing to reduce Gaza to rubble, killing, harming and destroying its people, and creating conditions of life calculated to bring about their physical destruction as a group".
It further alleges that"\textbf{acts} and \textbf{omissions} by Israel... are \textbf{genocidal} in character, as they are committed with the requisite specific \textbf{intent}... to destroy Palestinians in Gaza as a part of the broader Palestinian national, \textbf{racial} and \textbf{ethnical} group"and that"the conduct of Israel -- through its State organs, State agents, and other persons and entities \textbf{acting} on its instructions or under its direction, control or influence -- in relation to Palestinians in Gaza, is in \textbf{violation} of its \textbf{obligations} under the \textbf{Genocide} \textbf{Convention}."South \textbf{African} President Cyril Ramaphosa has compared Israel ' s policies in Gaza and the occupied West Bank with his country ' s past apartheid regime of \textbf{racial} segregation imposed by the white-minority rule that ended in 1994.
Several human rights organisations have said that Israeli policies towards Palestinians amount to apartheid.
South Africa ' s appeal includes @ @ @ @ @ @ @ @ @ @ interim orders for Israel to"immediately suspend its military operations in and against Gaza."There are several other countries that have referred to \textbf{genocide} committed by Israel in Gaza.
These countries include Algeria, Bolivia, Brazil, Colombia, Cuba, Iran, Palestine, Turkiye, Venezuela, Bangladesh, Egypt, Honduras, Iraq, Jordan, Libya, Malaysia, Namibia, Pakistan, and Syria.
But it is not yet clear whether they would become a party in the South \textbf{African} \textbf{case}.
It was expected that particularly Arab and Islamic countries would become a party to South Africa ' s petition, especially since the procedures followed by the ICJ allow it.
But barring a few, none of these countries have backed South Africa ' s appeal.
Among the Muslim countries only Turkiye, Jordan and Malaysia have supported the South \textbf{African} \textbf{case}.
The Jordanian government said that it ' s preparing a \textbf{legal} file to follow up on the \textbf{case}.
What is most perturbing is the dubious role played by Saudi Arabia, @ @ @ @ @ @ @ @ @ @ has not even publicly declared the Israeli action of killing of women and children in Gaza as \textbf{genocide}.
Saudi Arabia is also among the Muslim countries that had reportedly opposed any punitive action against Israel by the Organisation of Islamic Countries ( OIC ).
Although Pakistan has called Israel ' s actions in Gaza as \textbf{genocide}, sadly it has not yet given any indication of joining the South \textbf{African} petition.
Islamabad needs to play a much more active role in mobilising the \textbf{international} community to end the destruction of Gaza and killing of children by Israeli forces.
Some of the recent statements by the caretaker prime minister on Israel ' s war against Gaza have added to our policy confusion on the issue.
Meanwhile, Israel has mounted a concerted campaign to prevent the ICJ from concluding that it has indeed committed \textbf{genocide} in the Gaza Strip.
Last week the Israeli foreign ministry instructed its embassies to pressure politicians and diplomats in their host countries to make statements opposing South Africa ' s \textbf{case} at the ICJ.
Israel is also expected to challenge @ @ @ @ @ @ @ @ @ @ start arguing.
It ' s going to be a lengthy process of \textbf{legal} battle on the issue and there is no indication that it could prevent the \textbf{genocide} of Palestinians in the occupied territory.
But the South \textbf{African} \textbf{case} will help mobilise public opinion across the world against Israel ' s \textbf{genocidal} actions in Gaza and colonisation of Palestine.
The \textbf{international} community, particularly the Arab and Muslim countries, need to do more to stop the killing of women and children by Israeli forces before it ' s too late.

% matched lemmas: act, african, case, convention, court, crime, ethnical, genocidal, genocide, intent, international, legal, measure, obligation, omission, racial, resolution, submit, violation
\end{textsample}
